%% sections/discussion.tex
\section{Discussion, Limitations, and Future Work}
\label{sec:discussion}

\subsection{Key Findings}

This work establishes several findings relevant to low-resource AMR parsing more broadly.

\textbf{Silver annotations are a viable bootstrapping strategy.} The 97\% human acceptance rate of Gemini 2.5 Pro annotations demonstrates that a capable LLM can produce high-quality AMR annotations for a low-resource language when given well-structured prompts. The fine-tuned model trained on these silver annotations achieves a respectable F1\textsubscript{all} of 0.271, showing that silver-quality data is sufficient to bootstrap a working parser. This contrasts with the conventional assumption that only human expert annotations are reliable enough for training.

\textbf{Fine-tuning dramatically outperforms zero-shot transfer.} Even a small dataset of 880 training examples doubles the overall Smatch score (0.121 → 0.271) and halves the proportion of unparseable outputs (41.6\% → 18.2\%). This is a strong signal that the low-resource barrier is not insurmountable with modest annotation investment when starting from a multilingual pre-trained backbone.

\textbf{The language-adaptive pre-training did not help.} The most surprising finding is the near-total failure of System 3. We attribute this to a mismatch between the pre-training objective and the model's task expectations. In our implementation, the pre-training stage used \texttt{DataCollatorForSeq2Seq} with \texttt{labels = input\_ids}, effectively training the model to copy its input to its output as a reconstruction task. Unlike the original mBART denoising objective (which applies token masking, deletion, and sentence permutation before reconstruction), the copying objective places the encoder and decoder in a degenerate configuration: the encoder produces a rich contextual representation of the source, while the decoder learns to reproduce the same sequence verbatim. As evidence, the validation loss after pre-training converged to effectively zero (perplexity $\approx 1.0$), which is a signature of near-perfect memorisation rather than meaningful generalisation.

When this reconstructed model is subsequently fine-tuned for AMR, the decoder's strong prior towards copying the input conflicts with the AMR generation objective, which requires producing a structurally different output (a linearised graph) in response to a Konkani sentence. The result is a model that generates corrupted or ungrammatical linearisations for the vast majority of test sentences, causing the Penman parser to reject them. The 7 cases that \emph{do} parse (and score an average F1 of 0.437) may correspond to sentences where the copying prior accidentally aligned with generating something that resembles an AMR graph.

\subsection{Implications}

The failure of System 3 offers a practical lesson: language-adaptive pre-training for seq2seq models must use a proper denoising objective (token masking, span infilling, sentence shuffling) rather than reconstruction, to avoid the degenerate copy-mode that destroys the decoder's ability to generate novel output. Future work should implement the full mBART denoising collator—\texttt{DataCollatorForLanguageModeling} or a custom span-infilling collator—before pre-training on Konkani text.

\subsection{Limitations}

\textbf{Dataset size.} With 1,100 annotated sentences, KonkaniAMR is small compared to English AMR 3.0 (approximately 59,000 sentences) or even the mid-scale multilingual AMR datasets (XL-AMR has around 1,000 sentences per language). The limited training set means the fine-tuned parser has not seen many predicate types and is unlikely to generalise to rare or complex semantic phenomena.

\textbf{Silver annotation quality.} Although human verification produced a 97\% acceptance rate, the remaining 3\% required corrections, and the review was done by a single annotator, precluding inter-annotator agreement measurement. The quality of the silver annotations also depends on Gemini 2.5 Pro's ability to understand Konkani semantics, which has not been independently validated.

\textbf{Language code mismatch.} The use of \texttt{hi\_IN} as a proxy for Konkani in the mBART-50 tokeniser is an acknowledged approximation. Konkani and Hindi share many cognates but differ significantly in vocabulary, morphological patterns, and phonology. A tokeniser trained specifically on Konkani text would almost certainly improve subword segmentation quality and downstream performance.

\textbf{Evaluation metric limitations.} Smatch measures triple-level overlap and does not capture all aspects of parsing quality. It does not penalise wrong predicate senses if they have the same argument structure as the gold, and it does not directly measure the accuracy of named entity construction or coreference handling.

\textbf{Computational scale.} Our experiments were conducted on a single A100 80\,GB GPU. Larger-scale experiments—more epochs, larger batch sizes, hyperparameter search—were not feasible and may yield different conclusions.

\subsection{Future Directions}

Several directions emerge naturally from this work:

\begin{enumerate}
    \item \textbf{Proper denoising pre-training.} Reimplementing the pre-training stage with the mBART denoising objective (span masking + sentence permutation) is the highest-priority follow-up. We expect this to recover the benefit of Konkani-specific pre-training.

    \item \textbf{Dataset expansion.} Scaling KonkaniAMR to 5,000–10,000 sentences—potentially using other Konkani text sources (Wikipedia, court documents, literature) and the same LLM-annotation + human-verification pipeline—would enable training significantly stronger parsers.

    \item \textbf{Cross-lingual transfer from Indic languages.} Hindi and Marathi share substantial lexical overlap with Konkani. Pre-training or co-training on AMR data from these languages before fine-tuning on Konkani could provide additional signal.

    \item \textbf{Model-specific tokenisation.} Adding Konkani-specific vocabulary to the mBART-50 tokeniser—by training a sentencepiece model on our Konkani corpus and merging it—could reduce tokenisation noise.

    \item \textbf{Downstream applications.} Having established the first Konkani AMR parser, future work can explore its utility for information extraction from Goan news archives, question answering over Konkani encyclopaedic text, and machine translation evaluation.
\end{enumerate}

\subsection{Conclusion}

We have introduced KonkaniAMR, the first Abstract Meaning Representation dataset for the Konkani language, comprising 1,100 silver-annotated sentences verified by a human expert at a 97\% acceptance rate. We have established baselines for Konkani AMR parsing through zero-shot transfer, supervised fine-tuning, and language-adaptive pre-training followed by fine-tuning. Fine-tuning on KonkaniAMR more than doubles the zero-shot Smatch score and substantially reduces the rate of unparseable outputs, demonstrating that modest annotation investment can produce a functional low-resource AMR parser. Our analysis of the pre-training failure provides actionable guidance for future experiments. We release the dataset and all code to support continued research into semantic parsing for under-resourced South Asian languages.

\section*{Limitations}
\label{sec:limitations}

This work has the following limitations, discussed in more detail in Section~\ref{sec:discussion}: (i) the dataset is small (1,100 sentences), limiting generalisation; (ii) silver annotations produced by an LLM were verified by a single annotator without inter-annotator agreement measurement; (iii) the use of Hindi as a proxy language code for Konkani in the tokeniser is a known approximation; (iv) the language-adaptive pre-training used a reconstruction objective rather than a denoising objective, leading to degraded AMR generation; (v) results are reported on a single test split without statistical significance testing.
